\subsection{Two arrows in every hom-set around a middle corner}

\begin{theorem}[The double-hom middle-corner obstruction]
\label{C:middle-double-homs}
Let $C$ have objects $0,1,2$, with endomorphism monoids
$\{1\},M,\{1\}$ respectively, no backward arrows, and
\[
 |C(0,1)|=|C(1,2)|=|C(0,2)|=2.
\]
Here $M$ is any finite monoid.  Then $BC$ cannot have the
homotopy type of $\mathbb{CP}^2$.
More precisely, after relabeling the three two-point hom-sets,
the composition map has one of the following forms:
\[
\begin{array}{c|c}
\text{composition type}&\text{obstruction}\\\hline
\text{not onto}&H_1(BC;\mathbb Q)\ne0\\
\text{row or column}&\pi_1(BC)\cong\mathbb Z\\
\text{Latin square}&\pi_1(BC)\cong C_2\\
\text{one exceptional entry}&H_2(BC;\mathbb Q)=0.
\end{array}
\]
\end{theorem}
\begin{proof}
Write $S=C(0,1)$, $T=C(1,2)$ and $W=C(0,2)$, with
elements $s_i,t_i,w_i$ indexed by $i\in\{0,1\}$.
Composition is a map
\[
 \kappa:T\times S\longrightarrow W,
 \qquad (t,s)\longmapsto ts,
\]
balanced for the right $M$-action on $T$ and left $M$-action
on $S$.  If its image is a singleton, the unused arrow in
$W$ belongs to no nondegenerate simplex of dimension at least
two.  It is an additional edge between two vertices already
joined in the rest of the nerve, so it contributes a free
summand to first homology.

Suppose that $\kappa$ is onto.  A binary $2$ by $2$ matrix,
up to row, column and value interchanges, is a row matrix, a
column matrix, the Latin square $i+j\pmod2$, or the matrix
with its unique value $1$ at $(1,1)$.  These are exhaustive:
a value occurring once gives the last case, while two
occurrences are adjacent or diagonal.

For the row case, $\kappa(t_i,s_j)=w_i$.  Balance forces
every $M$-actor on $T$ to be the identity.  Choose $s_0,t_0$
as a maximal tree in the nerve's one-skeleton.  In its
fundamental-group presentation they become $1$, and the
composition relations give
\[
 w_0=1,\qquad s_1=1,\qquad w_1=t_1.
\]
The relations $t_0m=t_0$ kill every middle monoid generator.
No remaining relation involves the free generator $t_1$.
Thus the fundamental group is $\mathbb Z$.  The column case
is its opposite.

For the Latin square, balance forces a character
$c:M\to C_2$ such that both actions are translation by
$c(m)$.  With the same tree, composition gives
\[
 w_0=1,\qquad s_1=t_1=w_1,\qquad w_1^2=1,
 \qquad m=w_1^{c(m)}.
\]
The middle multiplication relations impose no further
relation because $c$ is a character.  Conversely these
assignments define a functor from $C$ to the one-object
groupoid $C_2$ and send $w_1$ to its nonidentity element.
The resulting fundamental group is exactly $C_2$.

It remains to handle
\[
 \kappa(t_i,s_j)=w_{ij},
\]
where $ij$ is ordinary multiplication of the binary indices.
If the row and column actions of $m$ are $\tau_m,\sigma_m$,
balance says
\[
 \tau_m(i)j=i\sigma_m(j)\qquad(i,j\in\{0,1\}).
\]
Putting one index equal to zero and then both equal to one
shows that the action pair is either $(\mathrm{id},\mathrm{id})$
or $(c_0,c_0)$, where $c_0$ is the constant map to zero.
There is consequently a character
$\chi:M\to\{0,1\}$ describing the common action on both
augmentation lines.  The following explicit chain
calculation proves the last assertion in the table.
\end{proof}

\begin{lemma}[The one-exception composition chain complex]
\label{C:AND-middle-bar}
In the last case of Theorem~\ref{C:middle-double-homs}, put
$R=\mathbb QM$ and let $\mathbb Q_\chi$ denote the
one-dimensional module given by $\chi$, on either side.
Then
\[
 H_1(BC;\mathbb Q)=H_2(BC;\mathbb Q)=0,
 \qquad
 H_n(BC;\mathbb Q)\cong
 \operatorname{Tor}_{n-2}^{R}
          (\mathbb Q_\chi,\mathbb Q_\chi)
 \quad(n\geq3).
\]
\end{lemma}
\begin{proof}
Let $C_0\subset C$ contain all objects, all middle
endomorphisms, and just the arrows $s_0,t_0,w_0$ between
different objects.  It is a subcategory because these states
are fixed by every middle actor and $t_0s_0=w_0$.
Its object $0$ is initial and its object $2$ terminal.
Thus $BC_0$ is contractible, and the normalized relative
nerve chains compute the reduced homology of $BC$.

Here is the full decomposition of those relative chains.
Write
\[
 e=s_1-s_0,\qquad f=t_1-t_0,\qquad v=w_1-w_0.
\]
As actual left and right modules,
\[
 \mathbb QS=\mathbb Qs_0\oplus\mathbb Q_\chi e,
 \qquad
 \mathbb QT=\mathbb Qt_0\oplus\mathbb Q_\chi f.
\]
For a right module $U$ and left module $V$, denote by
$B_*(U,R,V)$ the normalized two-sided monoid bar.  Its
degree-$q$ basis uses $q$ nonidentity middle arrows.  A
simplex with only an $S$-crossing contributes
$B_*(\mathbb Q,R,\mathbb Q_\chi)$ shifted by one; a
simplex with only a $T$-crossing contributes
$B_*(\mathbb Q_\chi,R,\mathbb Q)$ shifted by one.
The coefficient splittings just displayed make the
two-crossing simplices the direct sum of the three bars
with endpoint coefficients
\[
 (t_0,e),\qquad (f,s_0),\qquad (f,e),
\]
each shifted by two.  The fourth endpoint summand
$(t_0,s_0)$ lies in $C_0$ and is zero in relative chains.
Finally the long arrow contributes $\mathbb Qv$ in degree
one and in no other degree.

These are direct-sum decompositions of chain complexes,
with the following two endpoint differentials.  Deleting
the final crossing on the $(t_0,e)$ bar gives, with its
usual boundary sign, the identity onto the single-$S$
bar.  Its other endpoint deletion is zero since the
augmentation of $e$ is zero.  Likewise the $(f,s_0)$ bar
maps identically, up to sign, onto the single-$T$ bar and
has zero other endpoint deletion.  Their length-two
composition faces vanish because
\[
 \kappa(t_0,e)=0,\qquad \kappa(f,s_0)=0.
\]
Each of these two pairs is therefore the mapping cone of
the identity on its normalized bar, with a degree shift;
both are contractible by the identity-cone contraction.
All middle faces preserve the displayed summands, since
the module splittings are actual $R$-module splittings.
This verifies the cancellation on every simplex, rather
than only on a homology spectral sequence.

On the remaining $(f,e)$ bar both endpoint deletions are
zero.  Its only extra face is the degree-two composition
map
\[
 f\otimes e\longmapsto
 w_1-w_0-w_0+w_0=v.
\]
Consequently the remaining relative complex has
$B_{n-2}(\mathbb Q_\chi,R,\mathbb Q_\chi)$ in degree
$n\geq2$, $\mathbb Qv$ in degree one, and differential
from degree two to one given by this augmentation.
The degree-zero homology of that bar is
\[
 \mathbb Q_\chi\otimes_R\mathbb Q_\chi\cong\mathbb Q:
\]
the balancing relation on an element $m$ equates the
same scalar $\chi_m$ on the two sides.  The displayed
composition map induces the identity on this copy of
$\mathbb Q$.  It therefore kills both degree-one and
degree-two homology of the relative complex.  All higher
degrees are precisely the claimed shifted Tor groups.
\end{proof}

\begin{lemma}[The rational Latin-square calculation]
\label{C:Latin-middle-bar}
In the Latin-square case of Theorem~\ref{C:middle-double-homs},
let $\chi:M\to\{1,-1\}$ be its sign character.  Then
\[
 H_1(BC;\mathbb Q)=H_2(BC;\mathbb Q)=0,
 \qquad
 H_n(BC;\mathbb Q)\cong
 \operatorname{Tor}_{n-2}^{\mathbb QM}
       (\mathbb Q_\chi,\mathbb Q_\chi)
 \quad(n\geq3).
\]
In particular the obstruction survives a rational homology
and cohomology-ring equivalence to this category; the
fundamental-group computation alone is not needed for that use.
\end{lemma}
\begin{proof}
Use the rational vectors
\begin{gather*}
 s_+=\tfrac12(s_0+s_1),\quad
 t_+=\tfrac12(t_0+t_1),\quad
 w_+=\tfrac12(w_0+w_1),\\
 e=s_1-s_0,\qquad f=t_1-t_0,\qquad v=w_1-w_0.
\end{gather*}
The first two averaged vectors are fixed by every $M$-actor,
and the two augmentation lines have the same sign character.
The composition map satisfies
\[
 \kappa(t_+,s_+)=w_+,
 \quad\kappa(t_+,e)=\kappa(f,s_+)=0,
 \quad\kappa(f,e)=-2v.
\]

Collapse all three hom-sets to singletons, retaining $M$,
to obtain a category $C'$ with initial object $0$ and
terminal object $2$.  The induced normalized rational
chain map has a literal section: replace each singleton
cross arrow by its averaged vector, and on a simplex with
both adjacent crossings use their tensor product.  The
three identities just displayed, invariance of the averages,
and their augmentation one show that this section commutes
with every face and degeneracy.  Its image is therefore
the normalized chain complex of the contractible $BC'$.
This is a section of rational simplicial vector spaces;
no section by actual arrows is asserted.

Quotienting by this contractible summand leaves the same
two single-crossing bars and three double-crossing bars as
in Lemma~\ref{C:AND-middle-bar}, with $s_+,t_+$ replacing
$s_0,t_0$.  The two mixed pairs are again literal identity
cones.  The remaining complex has the two-sided
$\chi,\chi$ bar shifted by two and $\mathbb Qv$ in degree
one, with degree-two augmentation multiplication by $-2$.
Since
$\mathbb Q_\chi\otimes_{\mathbb QM}\mathbb Q_\chi\cong\mathbb Q$,
this is an isomorphism on its degree-zero bar homology.
The low-degree vanishing and higher shifted Tor formula
follow exactly as before.
\end{proof}

\begin{corollary}[Rational double-hom obstruction]
\label{C:middle-double-homs-rational}
Every category in Theorem~\ref{C:middle-double-homs} has
either nonzero first rational homology or zero second
rational homology.  Thus no space rationally homology
equivalent to such a category has the rational homology
of $\mathbb{CP}^2$.
\end{corollary}
\begin{proof}
Use the first-homology or fundamental-group computations
in the non-surjective, row and column cases, and
Lemmas~\ref{C:AND-middle-bar} and~\ref{C:Latin-middle-bar}
in the two remaining cases.
\end{proof}
